% Avsnitt 16.9, ur lectures/L16/README.md.

\section{Sammanfattning}
\label{c16:sec:review}

\needspace{6\baselineskip}
\subsection*{Det här ska du kunna nu}
\begin{itemize}
  \item Rita CAN-ramens tillståndsdiagram och säga vilket fält varje tillstånd svarar mot.
  \item Förklara varför varje fält behöver ett \code{LOAD}-tillstånd före sitt skifttillstånd.
  \item Implementera sändvägen: ladda skiftregistret, skifta ut, och gå vidare vid rätt puls.
  \item Driva bussen med öppen dränering korrekt: \code{bus_en} styr, \code{tx_bus} bestämmer vad.
  \item Grinda CRC-motorn rätt, och säga varför fel grindning ger en CRC som ser rimlig ut och är
    fel.
\end{itemize}

\needspace{6\baselineskip}
\subsection*{Frågor att testa dig själv med}
\begin{enumerate}
  \item Varför behöver varje fält ett eget \code{LOAD}-tillstånd?
  \item En nod som sänder släpper bussen under ACK-luckan. Varför, och vad förväntar den sig se där?
  \item \code{bus_en} och \code{tx_bus} är två separata signaler. Vad hade gått förlorat med bara
    en?
  \item Varför grindas \code{crc15} på \code{bit_valid and not stuff} och inte på \code{bit_done}?
  \item Vilka fält täcks av CRC:n, och var slutar täckningen exakt?
  \item Varför är tillståndsmaskinen en Mooremaskin här?
\end{enumerate}

\needspace{6\baselineskip}
\subsection*{Riktvärde}
Sändvägen bör vara på plats när \chapref{c17:ch} börjar, eftersom mottagarvägen byggs in i samma
maskin. Kommer du inte hela vägen: gör klart sändningen fram till och med CRC-fältet, och lämna ACK
och EOF till nästa pass.

\needspace{6\baselineskip}
\subsection*{Blicken framåt}
\Chapref{c17:ch} färdigställer \code{can_controller} med mottagning och arbitrering, och kör
systemtestbänken hela vägen.
